% Part: lambda-calculus
% Chapter: syntax
% Section: beta

\documentclass[../../../include/open-logic-section]{subfiles}

\begin{document}

\olfileid{lam}{syn}{bet}
\olsection{اختزال~$\beta$}

عندما نرى~$(\lambd[m][(\lambd[y][y]) m])$، فمن الطبيعي أن نخمن أن له
صلة ما بـ$\lambd[m][m]$، وهي أن الحد الثاني ينبغي أن يكون نتيجة
``تبسيط'' الأول. ويصوغ مفهوم \emph{اختزال~$\beta$} هذا الحدس صوريًا.

\begin{defn}[تقلّص~$\beta$، $\bredone$] \ollabel{defn:betacontr}
  \emph{تقلّص~$\beta$} ($\bredone$) هو أصغر علاقة متوافقة على الحدود
  تستوفي الشرط الآتي:
  \[
    (\lambd[x][N])Q \bredone \Subst{N}{Q}{x} 
  \]
  نقول إن~$P$ \emph{يتقلّص بحسب~$\beta$} إلى~$Q$ إذا كان~$P \bredone
  Q$. ويسمى الحد من الصورة~$(\lambd[x][N])Q$ \emph{موضع اختزال}.
\end{defn}

\begin{prob} \ollabel{prob:def}
  اكتب التعريفات الاستقرائية المكافئة لتقلّص~$\beta$ بالتفصيل، كما فعلنا
  لتغيير المتغير المقيد في~\olref[lam][syn][alp]{defn:aconvone}.
\end{prob}
  
\begin{defn}[اختزال~$\beta$، $\bred$] \ollabel{defn:betared}
  \emph{اختزال~$\beta$} ($\bred$) هو أصغر علاقة انعكاسية ومتعدية على
  الحدود تحتوي~$\bredone$. ونقول إن~$P$ يختزل بحسب~$\beta$ إلى
  $Q$ إذا كان~$P \bred Q$.
\end{defn}

سنكتب~$\redone$ بدل~$\bredone$، و$\red$ بدل~$\bred$، متى كان السياق
واضحًا.

وبعبارة غير رسمية، يكون~$M \bred N$ إذا وفقط إذا أمكن تغيير~$M$ إلى
$N$ بصفر أو بعدة خطوات من تقلّص~$\beta$.

\begin{defn}[السواء بحسب~$\beta$]
يقال إن الحد الذي لا يمكن تقليصه بحسب~$\beta$ أكثر \emph{سوي بحسب~$\beta$}.
\end{defn}

إذا كان~$M \bred N$ وكان~$N$ سويًا بحسب~$\beta$، قلنا إن~$N$
\emph{صورة سوية} لـ$M$. وقد يسأل المرء هل الصورة السوية لحد وحيدة؛
والجواب نعم، كما سنرى لاحقًا.

لنتأمل بعض الأمثلة.
\begin{enumerate}
\item لدينا
\begin{align*}
(\lambd[x][xxy]) \lambd[z][z] & \redone (\lambd[z][z])(\lambd[z][z]) y \\
& \redone (\lambd[z][z]) y \\
& \redone y
\end{align*}
\item قد يجعل ``تبسيط'' حد ما ذلك الحد أشد تعقيدًا:
\begin{align*}
(\lambd[x][xxy])(\lambd[x][xxy]) & \redone (\lambd[x][xxy])(\lambd[x][xxy])y \\
& \redone (\lambd[x][xxy])(\lambd[x][xxy])yy \\
& \redone \dots
\end{align*}
\item وقد يترك حدًا من دون تغيير أيضًا:
\[
(\lambd[x][xx])(\lambd[x][xx]) \redone (\lambd[x][xx])(\lambd[x][xx])
\]
\item ويمكن كذلك اختزال بعض الحدود بأكثر من طريقة؛ فمثلًا،
\[
(\lambd[x][(\lambd[y][yx]) z]) v \redone (\lambd[y][yv]) z
\]
بتقليص التطبيق الأبعد إلى الخارج؛ و
\[
(\lambd[x][(\lambd[y][yx]) z]) v \redone (\lambd[x][zx]) v
\]
بتقليص التطبيق الأبعد إلى الداخل. لكن لاحظ أن كلا الحدين يختزل بعد ذلك،
في هذه الحالة، إلى الحد نفسه~$zv$.
\end{enumerate}

ليست النتيجة النهائية في المثال الأخير مصادفة، بل توضح خاصية عميقة
ومهمة لحساب لامبدا تعرف باسم خاصية تشيرش--روسر.

\begin{digress}
  توجد عمومًا أكثر من طريقة لاختزال حد بحسب~$\beta$، ولذلك ابتكرت
  استراتيجيات اختزال كثيرة، أشهرها \emph{الاستراتيجية الطبيعية}. وتقلّص
  الاستراتيجية الطبيعية دائمًا موضع الاختزال \emph{الأبعد إلى اليسار}،
  حيث يعرّف موضع الاختزال بنقطة بدايته في الحد. وللاستراتيجية الطبيعية
  الخاصية المفيدة الآتية: يمكن اختزال حد إلى صورة سوية باستراتيجية ما إذا
  وفقط إذا أمكن اختزاله إلى صورة سوية باستعمال الاستراتيجية الطبيعية.
  وسنستعمل فيما يأتي الاستراتيجية الطبيعية ما لم ينص على غير ذلك.
\end{digress}

\begin{defn}[التكافؤ بحسب~$\beta$، $\equal$]
  \emph{التكافؤ بحسب~$\beta$} ($\equal$) هو العلاقة المعرفة استقرائيًا
  كما يأتي:
  \begin{enumerate}
  \item $M \equal M$.
  \item إذا كان~$M \equal N$، فإن~$N \equal M$.
  \item إذا كان~$M \equal N$ و$N \equal O$، فإن~$M \equal O$.
  \item إذا كان~$M \equal N$، فإن~$PM \equal PN$.
  \item إذا كان~$M \equal N$، فإن~$MQ \equal NQ$.
  \item إذا كان~$M \equal N$، فإن~$\lambd[x][M] \equal \lambd[x][N]$.
  \item $(\lambd[x][N])Q \equal \Subst{N}{Q}{x}$.
  \end{enumerate}
\end{defn}

تجعل القواعد الثلاث الأولى العلاقة علاقة تكافؤ؛ وتجعلها القواعد الثلاث
التالية متوافقة؛ وتضمن القاعدة الأخيرة أنها تحتوي تقلّص~$\beta$.

وبعبارة غير رسمية، يكون حدان متكافئين بحسب~$\beta$ إذا وفقط إذا أمكن
تغيير أحدهما إلى الآخر بصفر أو بعدة خطوات من تقلّص~$\beta$ أو من
``معكوس'' تقلّص~$\beta$. ويعرّف معكوس تقلّص~$\beta$ بحيث يتقلّص~$M$
عكسيًا بحسب~$\beta$ إلى~$N$ إذا وفقط إذا تقلّص~$N$ بحسب~$\beta$ إلى
$M$.

وسنوسع العلاقة، إلى جانب القواعد السابقة، بقواعد أخرى، ونرمز إلى علاقة
التكافؤ الموسعة بـ$\equal[X]$، حيث~$X$ هي القاعدة الموسعة.

\end{document}
